\documentclass[runningheads,a4paper]{llncs}

\usepackage{amsmath,amssymb}
\usepackage{graphicx}
\usepackage{booktabs}
\usepackage{multirow}
\usepackage{siunitx}
\usepackage{url}
\usepackage{hyperref}
\usepackage{xcolor}

\hypersetup{
    pdftitle={Energy-Aware Priority Scheduling in Smart Sensor Networks Using Policy-Guided Deep Reinforcement Learning},
    pdfauthor={Anand Emmanuel Raju B., Radhika N., Radhika G.},
    pdfsubject={ICTSES 2026 Camera-Ready Manuscript},
    pdfkeywords={Smart Energy Systems, Wireless Sensor Networks, Reinforcement Learning, ns-3, QoS, Priority Scheduling}
}

\title{Energy-Aware Priority Scheduling in Smart Sensor Networks Using Policy-Guided Deep Reinforcement Learning}
\titlerunning{Energy-Aware Priority Scheduling in Smart Sensor Networks}

\author{Anand Emmanuel Raju B.\inst{1} \and Radhika N.\inst{1} \and Radhika G.\inst{1}}
\authorrunning{A. E. R. Billa et al.}

\institute{
Department of Computer Science and Engineering,\\
Amrita School of Computing, Coimbatore,\\
Amrita Vishwa Vidyapeetham, India\\
\email{cb.sc.p2cse25003@cb.students.amrita.edu, n\_radhika@cb.amrita.edu, g\_radhika@cb.amrita.edu}
}

\begin{document}

\maketitle

\begin{abstract}
Smart Grid sensor networks have a basic problem: ensur-
ing reliable delivery of protection critical information (fault notifications
and breaker commands) to the operators while maintaining continu-
ous telemetry channels (voltage, current, temperature). Scheduling with
Strict Priority policy—which always transmits alarms first—guarantees
better latency performance for the critical data but prevents telemetry
from working during disturbances, which is crucial at those times. In
this paper, we explore how reinforcement learning can solve competing
scheduling priorities problem better than heuristic-based approaches. A
Double-DQN policy was trained in the environment created by ns-3 sim-
ulator and PyTorch and evaluated on 20 matched runs. We show that
there is an obvious trade-off: critical packets transmit 5.2 faster with
Strict Priority policy (0.0274 seconds vs. 0.1431 seconds), but telemetry
lags about twice longer (0.6171 seconds vs. 0.4287 seconds). The pro-
posed learned policy provides a latency improvement of 30.5\% at a cost
of 2.56\% degradation of packet delivery ratio.

\keywords{Smart Energy Systems \and Wireless Sensor Networks \and Reinforcement Learning \and ns-3 \and QoS \and Priority Scheduling}
\end{abstract}

\section{Introduction}
At the distribution side, the power network highly relies on wireless sensor net-
works to carry out the following tasks: detection of faulty equipment, monitoring
the condition of the network in real time, and remote operation of breakers. In re-
ality, each gateway sensor is required to concurrently process two different types
of data traffic. The first one is critical in nature, comprising protection alarms
that require millisecond delivery. The second one is non-critical, and comprises
telemetry (voltage, current, temperature).

The systems employed today employ Strict Priority scheduling in handling
this issue. The simple principle behind the scheduling is that critical packets
should always come first while routine telemetry should be sent only if there are
no packets in the critical queue. It works quite well in the case of critical alarms
as the operators do get their alarms. However, in the case of a power grid failure,
where the protection alarms occur heavily and continuously for a few seconds,
the routine telemetry queue gets totally starved out.

Is there a way machine learning could do better? The answer seems quite
obvious: rather than following a hardcoded rule, train the gateway to come up
with an adaptive policy that would balance the two traffic goals. It is possible
that the learning system will realize that the recently received critical message
should not be forwarded immediately since its deadline is still 500 milliseconds
ahead of time, that is the moment when it could flush out old telemetry packets
before they become too outdated. However, this is not an easy task to accomplish.

Our study explicitly attempts to solve this problem. We used a Double DQN
algorithm for scheduling with explicit reward shaping to prevent collapse, trained it
in ns-3 simulation combined with PyTorch, and conducted 20 matched experi-
ments against Strict Priority and Random algorithms. The important obser-
vation that we made is straightforward: our algorithm does not surpass Strict
Priority in terms of alert latency. Critical packets are delivered in 0.1431 sec-
onds via our algorithm, compared to 0.0274 seconds via Strict Priority. However,
there is a trade-off here. Our scheduler delivers telemetry data with 30.5%
lower latency (0.4287 seconds compared to 0.6171 seconds).

The core contributions include:
\begin{itemize}
    \item Fully reproducible simulation framework that integrates ns-3.40 and Py-
Torch for a direct comparison of policies in a controllable environment. The
framework is open source and extendable.
    \item Guidance based on a soft reward system that avoids policy breakdowns, while
preserving flexibility towards the attainment of dual objectives. Intensity of
guidance depends on the age of the packets, being less intensive if deadlines
are far and more intensive near expiration.
    \item Empirical results proving that the learned strategy generalizes well for three
network sizes (50, 100, 200 nodes), although success varies according to the
situation.
    \item Tradeoffs Quantified: 2.56\% reduction in mission-critical delivery (98.34\%
to 95.78\%) in exchange for 30.5\% decrease in telemetry delay. The ratio is
skewed towards operators who require awareness during grid disturbance
conditions.
\end{itemize}

\section{Related Work}
Application of reinforcement learning to wireless communications has been made
within the last decade. This is due to the simple reason that instead of manual
tuning of communication protocols, one would be able to adjust to changing
circumstances using learning-based approaches. Surveys \cite{frikha2021wrlsurvey,luong2019drlsurvey} on these topics con-
sider spectrum allocation, routing, MAC protocol design, and packet scheduling.
Nonetheless, both surveys highlight the same problem - the lack of guarantee of
safety in the process of RL and large amounts of data needed for training.

For the application of sensor networks, there are a number of implementation
and analysis approaches that touch on the issue. According to Leong \cite{leong2020automatica}, neural
policies are able to adapt to the distribution of packets arriving in the network
and fading behavior which is difficult for any fixed heuristics to capture. The
application of distributed RL in order to achieve coverage and lifetime has been
considered by Sharma and Chauhan \cite{sharma2020wsn}.

However, there is one area where a gap exists within the literature, namely
a direct comparison between the two types of scheduler in question. While most
tests of learning compare learning against random allocation and no control at
all, we will be filling that gap through comparing the two scheduling strategies,
with an eye towards understanding the tradeoffs between the two.

\section{System Model and Problem Formulation}

\subsection{Network and Traffic Model}
A sensing cluster that represents an edge device of a smart grid distribution
substation is developed. The cluster consists of 100 wireless sensors randomly
distributed over an area of 150 m × 150 m and communicating with a central
gateway. The physical and MAC layers employ ns-3.40 implementation of IEEE
802.15.4 in the 2.4 GHz band.

Sensor nodes generate two distinct traffic classes: High-Priority (HP) packets rep-
resenting critical protection alerts and fault notifications, and NormalPriority (NP)
packets representing telemetry and state-estimation feeds. HP generation follows
an unannounced bursty Poisson process; NP traffic is periodic. Each node main-
tains separate FIFO buffers (HP capacity: 50 packets, NP capacity: 100). Packets
exceeding deadlines ($T_H^{max} = 1.0$ s, $T_N^{max} = 3.0$ s) are dropped. The
gateway scheduler operates at 100 ms decision intervals over 5000-step episodes
(500 s real time).

\subsection{MDP Definition}
Scheduling is formulated as a discounted infinite-horizon MDP.

\begin{equation}
s_t = \left[q_t^N,\; q_t^H,\; e_t,\; a_t^N,\; a_t^H,\; d_t\right],
\end{equation}
where $q_t^N, q_t^H \in [0,1]$ denote normalized buffer occupancy; $e_t$ denotes
normalized residual energy; $a_t^N, a_t^H \in [0,1]$ denote normalized head-of-line
packet ages relative to expiration; and $d_t$ denotes normalized channel contention.

The action space comprises five discrete choices:
\begin{equation}
\mathcal{A} = \{\text{send HP (0)},\; \text{send normal (1)},\; \text{drop normal (2)},\; \text{sleep (3)},\; \text{drop HP (4)}\}.
\end{equation}
Actions 0, 1 attempt transmission; Action 3 enters sleep; Actions 2, 4 enable
buffer purging of aged packets.

\subsection{Reward Formulation with Soft Policy Guidance}
The reward mechanism faces the following challenge: the reward mechanism
based on simple throughput produces policy collapse (failure to take into account
important traffic), whereas punishment makes it impossible to avoid replicating
the Strict Priority approach. Solution:
\begin{equation}
r_t = \tanh\left(\frac{r_t^{base} + r_t^{guide}}{100}\right),
\end{equation}
where the base reward takes into account multiple criteria (successful deliveries
with weights depending on classes $w_1 = 25.0$ for HP, $w_2 = 5.0$ for NP;
penalization for packet drops and timeouts; penalization for latency; energy
consumption). The guidance reward provides reinforcement for HP delivery and
negative feedback for late HP packets only if urgent ($a_t^H > \tau_u = 0.55$),
leaving flexibility for trading off between the two. The base part takes into
account the main network metrics:
\begin{equation}
r_t^{base} = w_1 N_{sent}^H + w_2 N_{sent}^N + w_3 N_{drop}^H + w_4 N_{drop}^N + w_5 N_{to}^H + w_6 N_{to}^N + w_7 (a_t^H)^2 + w_8 (a_t^N)^2 - w_9 \Delta e_t,
\end{equation}
where $N_{sent}^H$ and $N_{sent}^N$ represent confirmed deliveries ($w_1 = 25.0, w_2 = 5.0$);
$N_{drop}^H, N_{drop}^N, N_{to}^H, N_{to}^N$ penalize buffer overflow drops and timeout
expirations ($w_3 = -40.0, w_4 = -8.0, w_5 = -80.0, w_6 = -10.0$); quadratic
terms penalize waiting queue latencies ($w_7 = -40.0, w_8 = -2.0$); and $\Delta e_t$
penalizes energy depletion ($w_9 = 250.0$). If the task is chosen to be sleep
(Action 3), the scheduler gets a conservation reward $r_{sleep} = +0.5$.

In order to guarantee the agent services important traffic without losing its
operational flexibility, we propose the soft guidance policy term:
\begin{equation}
r_t^{guide} =
\begin{cases}
+\alpha, & a_t = 0 \text{ and } q_t^H > 0\\
-\beta, & a_t \neq 0,\; q_t^H > 0,\; a_t^H > \tau_u\\
0, & \text{otherwise}
\end{cases}
\end{equation}
with bonus $\alpha = 6.0$, penalty $\beta = 8.0$, and urgency age threshold
$\tau_u = 0.55$. Under such a formulation, the agent gets positive rewards for
delivering the HP packets. Importantly, the penalty $\beta$ for not doing the
task is not levied when the waiting HP packet is young ($a_t^H \le \tau_u$).
This design grants the scheduler autonomy to temporarily withhold a newly
arrived HP packet to flush an aging normal queue, intervening with negative
feedback only when an HP packet approaches critical delay boundaries.

\section{Methodology}

\subsection{Learning Architecture}
We implement Double Deep Q-Network (Double-DQN) to mitigate maximization bias in
standard Q-learning. The target computation splits action selection, through
online network $\theta_t$, from evaluation, using target network $\theta_t^{-}$:
\begin{equation}
Y_t^{DoubleQ} = r_{t+1} + \gamma Q\left(s_{t+1}, \arg\max_a Q(s_{t+1}, a; \theta_t); \theta_t^{-}\right).
\end{equation}

Prioritized Experience Replay (PER) speeds up learning on rare critical
events using transition probabilities $P(i) = p_i^\mu / \sum_k p_k^\mu$ with
$p_i = |\delta_i| + \epsilon_{per}$ (temporal difference error) and $\mu = 0.6$.
Importance sampling weights $w_i = (N \cdot P(i))^{-\nu}$ account for bias,
$\nu = 0.4 \rightarrow 1.0$.

The network is a fully connected MLP: input layer (6 neurons), two hidden
layers (64 neurons, ReLU), output layer (5 Q-values). Target weights update via
soft Polyak averaging with $\tau = 0.005$:
\begin{equation}
\theta^- \leftarrow \tau \theta + (1-\tau)\theta^-.
\end{equation}
Optimization uses Adam ($lr = 1 \times 10^{-4}$, Cosine Annealing), gradient
clipping (max $L_2 = 1.0$), and $\epsilon$-greedy exploration ($1.0 \rightarrow 0.05$
decay). Evaluation uses $\epsilon = 0.0$.

\subsection{Baseline Schedulers}
The comparisons use two different baselines: (i) Strict Priority (SP): sends
HP if available, otherwise NP; (ii) Random: randomly selects actions.

\subsection{Training and Validation Protocol}
The training is done for 50 epochs with 5000 iterations per epoch. The checkpoint
selection criterion is based on the validation metric on fixed seeds:
\begin{equation}
S_{val} = 0.40 \left(\frac{PDR}{100}\right) + 0.10 \left(1 - \frac{PLR}{100}\right) + 0.35 \left(\frac{HR}{100}\right) + 0.15 \left(1 - \min\left(\frac{D_H}{1.0}, 1.0\right)\right),
\end{equation}
The early stopping criteria are activated following 15 epochs without improve-
ment. The evaluation process conducts 20 concurrent tests for each strategy
(seeds 12545–12564).

\begin{table}[t]
\centering
\caption{Simulation and Training Hyperparameters (V76 Configuration).}
\label{tab:setup}
\scriptsize
\setlength{\tabcolsep}{6pt}
\begin{tabular}{ll}
\toprule
\textbf{Parameter Category} & \textbf{Configured Value} \\
\midrule
Network Simulator & ns-3.40 interfaced via ns3-gym \\
Network Topology & 100 Sensor Nodes + 1 Central Sink (150\,m $\times$ 150\,m) \\
Physical / MAC Layer & IEEE 802.15.4 (2.4\,GHz ISM Band, CSMA/CA) \\
Control Interval ($\Delta t$) & 100\,ms per decision slot \\
Episode Duration & 5000 decision steps (500\,s real time) \\
Traffic Generation & Bursty Poisson (HP Alerts) + Periodic Telemetry (NP) \\
Queue Capacities & HP Buffer: 50 packets; Normal Buffer: 100 packets \\
State / Action Dimensionality & $|\mathcal{S}| = 6$ continuous, $|\mathcal{A}| = 5$ discrete \\
Neural Architecture & MLP: $6 \to 64 \to 64 \to 5$ (ReLU activations) \\
Replay Buffer Capacity & 100,000 transitions (Prioritized: $\mu=0.6, \nu=0.4 \to 1.0$) \\
Mini-batch Size & 128 transitions \\
Discount Factor ($\gamma$) & 0.99 \\
Optimizer / Learning Rate & Adam ($\text{lr} = 1 \times 10^{-4}$, Cosine Annealing) \\
Target Network Update & Soft Polyak update ($\tau = 0.005$) \\
Exploration Schedule & $\epsilon$-greedy: $1.0 \to 0.05$ (decay factor $0.99995$) \\
Evaluation Repetitions & 20 matched runs per policy (seeds $12545$--$12564$) \\
\bottomrule
\end{tabular}
\end{table}

\subsection{Statistical Evaluation Protocol}
The metrics used include throughput (kbps), packet delivery ratio (PDR\%),
packet loss ratio (PLR\%), average delays, energy consumption per bit, and critical packet delivery ratio. All are reported as mean ± 95\% confidence interval (Student’s t).

For pairwise significance testing, we verify normality of paired differences via Shapiro-Wilk test $\alpha = 0.05$. If normal, paired t-test is applied; otherwise Wilcoxon signed-rank test. Effect sizes are quantified using Cohen’s $d_z = \frac{D}{\sigma_D}$across the 20 paired runs.

\section{Results}
\label{sec:results}
The results show that the proposed DQN policy significantly improves normal-
traffic delay while preserving a high HP delivery ratio.

\begin{table}[t]
\centering
\caption{WSN Performance Metrics Across 20 Matched Evaluation Runs (Mean $\pm$ 95\% CI).}
\label{tab:main_results}
\resizebox{\textwidth}{!}{%
\begin{tabular}{lccccccc}
\toprule
\textbf{Policy} & \textbf{Throughput} & \textbf{PDR} & \textbf{PLR} & \textbf{Normal Delay} & \textbf{HP Delay} & \textbf{Energy} & \textbf{HP Delivery} \\
 & \textbf{(kbps)} & \textbf{(\%)} & \textbf{(\%)} & \textbf{(s)} & \textbf{(s)} & \textbf{(nJ/bit)} & \textbf{Ratio (\%)} \\
\midrule
Proposed DQN-Edge & $7.0267 \pm 0.0702$ & $50.3291 \pm 0.6022$ & $40.3309 \pm 0.8373$ & \textbf{0.4287 $\pm$ 0.0261} & $0.1431 \pm 0.0026$ & $443.8708 \pm 16.9422$ & $95.7779 \pm 0.2916$ \\
Strict Priority (SP) & \textbf{7.9104 $\pm$ 0.0054} & \textbf{56.6625 $\pm$ 0.3137} & \textbf{34.7224 $\pm$ 0.5326} & $0.6171 \pm 0.0384$ & \textbf{0.0274 $\pm$ 0.0015} & \textbf{417.0370 $\pm$ 17.2341} & \textbf{98.3374 $\pm$ 0.1941} \\
Random Baseline & $2.8000 \pm 0.0350$ & $20.0709 \pm 0.3314$ & $71.4943 \pm 0.5450$ & $0.8545 \pm 0.0597$ & $0.3521 \pm 0.0129$ & $806.7023 \pm 22.7489$ & $34.4862 \pm 0.8563$ \\
\bottomrule
\end{tabular}%
}
\end{table}

\begin{table}[ht]
\centering
\caption{Matched Pairwise Statistical Comparison: Proposed DQN-Edge vs. Strict Priority ($n=20$ runs).}
\label{tab:stats}
\scriptsize
\setlength{\tabcolsep}{8pt}
\begin{tabular}{llcc}
\toprule
\textbf{Evaluated Metric} & \textbf{Statistical Test} & \textbf{\textit{p}-value} & \textbf{Paired Cohen's $d_z$} \\
\midrule
Throughput (kbps) & Paired $t$-test & $2.656 \times 10^{-16}$ & $-5.8046$ \\
PDR (\%) & Paired $t$-test & $1.655 \times 10^{-15}$ & $-5.2568$ \\
PLR (\%) & Paired $t$-test & $4.749 \times 10^{-13}$ & $+3.8504$ \\
Avg. Normal Delay (s) & Wilcoxon signed-rank & $1.907 \times 10^{-6}$ & $-3.3284$ \\
Avg. HP Delay (s) & Paired $t$-test & $3.038 \times 10^{-26}$ & $+19.6001$ \\
Energy (nJ/bit) & Paired $t$-test & $1.687 \times 10^{-14}$ & $+4.6310$ \\
HP Delivery Ratio (\%) & Paired $t$-test & $2.176 \times 10^{-11}$ & $-3.1005$ \\
\bottomrule
\end{tabular}
\end{table}

\subsection{Convergence and Training Stability}
Figure~\ref{fig:learning} illustrates the training progression over 30 elapsed
epochs prior to early stopping. The validation score trajectory shows rapid initial
policy improvement over the first 5 epochs, transitioning into stable, monotonic
refinement without policy collapse or loss divergence. The exploration rate
$\epsilon$ decays steadily, confirming that the agent transitions from exploratory
state-space coverage to stable deterministic execution. The model achieving the
highest validation score at Epoch 12 was preserved for full benchmarking.

\begin{figure}[htbp]
    \centering
    \includegraphics[width=0.72\columnwidth]{figures/fig_1_learning_curve.png}
    \vspace{-2mm}
    \caption{Double-DQN training progression and validation-score trend over 30 epochs (V76).}
    \label{fig:learning}
\end{figure}

\subsection{Comparative Analysis Against Heuristic Baselines}
The empirical results establish that the policy-guided DQN scheduler substan-
tially outperforms the unguided Random baseline across all performance metrics.
Compared to Random, DQN increases aggregate throughput from 2.8000 kbps to
7.0267 kbps (a 150.9\% gain), elevates overall PDR from 20.07\% to 50.33\%,
lowers average HP latency from 0.3521 s to 0.1431 s (a 59.4\% reduction), and
improves energy efficiency from 806.70 nJ/bit to 443.87 nJ/bit (a 45.0\%
energy reduction). Furthermore, the HP delivery ratio surges from 34.49\%
under random selection to 95.78\% under the learned controller.

When compared against the dedicated Strict Priority baseline, however, the
learned agent does not achieve superior throughput or HP delay. Strict Priority
delivers an average throughput of 7.9104 kbps, an overall PDR of 56.66\%, and
an ultra-low HP latency of 0.0274 s alongside an HP delivery ratio of 98.34\%.
The paired statistical tests in Table~\ref{tab:stats} confirm that these margins
are statistically significant ($p < 10^{-11}$), exhibiting large effect sizes.

\subsection{Delay Distribution and Buffer Dynamics}
Figure~\ref{fig:eval_plots} presents a visual comparison of means and 95\%
confidence intervals across key QoS indicators alongside the empirical Cumulative
Distribution Functions (CDFs) of packet latencies for both traffic classes.

\begin{figure}[htbp]
    \centering
    \begin{minipage}[b]{0.48\textwidth}
        \centering
        \includegraphics[width=\textwidth]{figures/fig_4_bar_charts.png}
        \caption{Performance metrics (Mean $\pm$ 95\% CI).}
        \label{fig:bar}
    \end{minipage}
    \hfill
    \begin{minipage}[b]{0.48\textwidth}
        \centering
        \includegraphics[width=\textwidth]{figures/fig_5_delay_cdf.png}
        \caption{Delay CDFs for traffic classes.}
        \label{fig:cdf}
    \end{minipage}
    \label{fig:eval_plots}
    \vspace{-2mm}
\end{figure}

The latency CDF curves in Figure~\ref{fig:cdf} reveal the underlying mechanics
of each scheduling policy. Strict Priority drives the HP delay curve aggressively
toward the origin, ensuring that over 90\% of all high-priority packets are
cleared within 50 ms. However, this aggressive service discipline severely de-
grades normal telemetry traffic; the normal delay distribution for Strict Priority
exhibits a long tail, with average latency climbing to 0.6171 s. In contrast,
the proposed DQN agent balances queue service: it accepts a modest shift in HP
delay (averaging 0.1431 s) to pull the normal traffic delay curve substantially
to the left. As reflected in Table~\ref{tab:main_results}, DQN reduces average
normal packet delay to 0.4287 s—a 30.5\% latency reduction compared to Strict
Priority ($p = 1.907 \times 10^{-6}, d_z = -3.3284$).

\subsection{Trade-Off Analysis in Smart Grid Telemetry}
The quantitative divergence between DQN-Edge and Strict Priority demonstrates a
clear engineering trade-off. Strict Priority behaves as a greedy, myopic heuris-
tic; it prioritizes immediate HP transmissions regardless of normal queue status.
In smart utility applications characterized by bursty fault activity, this myopic
behavior induces buffer overflow for routine measurements, depriving grid oper-
ators of continuous current, voltage, and thermal sensor readings.

The policy-guided DQN scheduler learns a cooperative balancing policy. By ob-
serving head-of-line packet ages ($a_t^H, a_t^N$), the agent learns that newly
arrived HP packets can tolerate brief queuing periods without violating their
1.0 s deadline. During these brief windows, the agent opportunistically clears
pending normal telemetry packets before they expire. By sacrificing a minor
2.56\% in high-priority delivery ratio and accepting a small latency increase of
115.7 ms, the DRL scheduler improves normal-priority packet latency by 30.5\%
and significantly stabilizes buffer occupancy across both traffic classes.

\subsection{Diagnostic Analysis: Reward Shaping vs. Policy Collapse}
Our development process provided an important lesson regarding reward formula-
tion in communications DRL. In earlier experiments using an unconstrained reward
function (V74 configuration), the scheduler suffered from complete policy col-
lapse: the agent converged to a degenerate policy that selected Action 1 (send
normal) in 100\% of decision steps. Because normal packets were abundant,
transmitting normal traffic provided steady rewards, whereas attempting HP
transmissions incurred risk due to bursty arrivals and contention. In that failed
run, zero HP packets were delivered. However, because the simulator contained a
zero-division safeguard defaulting average delay to 0.0 when no packets were de-
livered, the output superficially reported ``0.0000 s delay''—a deceptive artifact
that masked complete high-priority starvation.

Introducing the soft policy guidance in V76, paired with an HP-aware validation
scoring metric, successfully resolved this pathology. By penalizing the scheduler
only when an HP packet aged past $\tau_u = 0.55$, the agent was steered away
from trivial local optima without sacrificing its ability to dynamically manage
queues. This highlights that reinforcement learning applied to communication
protocols requires explicit per-class packet tracking, action distribution validation,
and tail latency profiling rather than aggregate reward monitoring alone.

\section{Discussion}
The results confirm that Double-DQN can learn stable, multi-objective schedul-
ing behaviors that enhance fairness in mixed-criticality sensor clusters. Never-
theless, the fact that Strict Priority remains superior on raw high-priority latency
indicates structural limitations in the unconstrained DRL formulation. First, the
compact 6-dimensional state vector may limit the agent's observability, hiding
fine-grained channel state variations. Second, single-agent Q-learning struggles
to match the extreme responsiveness of deterministic heuristics on boundary cases
without explicit safety overrides.

To advance learning-based scheduling for smart grid deployments, future work
should address three methodological directions:
\begin{itemize}
    \item \textbf{Constrained and Safe RL:} Formulating the scheduling task as a Constrained Markov Decision Process (CMDP) solved via Lagrangian relaxation or action projection layers to guarantee strict bounds on HP latency.
    \item \textbf{Distributional Value Approximation:} Utilizing Distributional RL (such as QR-DQN or IQN) to model the full distribution of queuing delays rather than expected values, mitigating tail-latency outliers.
    \item \textbf{Out-of-Distribution Sensitivity Analysis:} Testing trained policies against variable node densities, shifting HP-to-normal traffic ratios, and channel interference profiles to evaluate generalization beyond fixed training topologies.
\end{itemize}

\section{Threats to Validity}
Several factors may influence the generalizability of our empirical findings. First,
all evaluations were conducted within discrete-event simulation in ns-3.40. While
ns-3 provides high-fidelity modeling of 802.15.4 CSMA/CA protocols, physical
hardware deployments introduce non-idealities such as transceiver switching jitter,
clock drift, and complex spatial multipath fading that simulation abstractions omit.

Second, the observed performance boundaries are tied to the evaluated traffic
load (100 nodes, 100 ms intervals, specific deadline limits). Varying the ratio of
high-priority to normal traffic or increasing node contention could shift the trade-
off surface between Strict Priority and DRL. Finally, our experimental setup
focuses on single-hop cluster-head scheduling; extending these controllers to
multi-hop meshed forwarding topologies represents an open research challenge.

\section{Conclusion}
This study presented a reproducible, statistically grounded evaluation of deep
reinforcement learning for priority-aware packet scheduling in wireless sensor
networks supporting smart energy infrastructure. By coupling ns-3 with
PyTorch via ns3-gym, we trained a policy-guided Double-DQN agent over
mixed-criticality traffic workloads. Across 20 seed-synchronized evaluation runs,
the learned scheduler substantially outperformed random baseline allocation and
preserved a 95.78\% high-priority delivery ratio, but did not match Strict Priority
in raw high-priority latency or throughput. However, the learned policy
demonstrated superior multi-objective queue balancing, achieving a 30.5\%
reduction in normal telemetry delay and mitigating buffer starvation. These
findings provide an open empirical benchmark that clarifies the practical
performance boundaries and trade-offs of learning-based controllers in safety-
critical sensor communications.

\subsubsection*{Reproducibility Package.}
To facilitate independent verification, the complete simulation artifacts—including
the raw evaluation results (\texttt{raw\_results\_v76.csv}), training history
(\texttt{training\_log\_v76.csv}), statistical summary tables
(\texttt{final\_results\_v76.xlsx}), and the underlying ns-3 C++ and PyTorch source
code—are made available for academic reuse.

\subsubsection*{Disclosure of Interests.}
The authors have no competing interests to declare that are relevant to the content of this article.

\nocite{*}
\bibliographystyle{splncs04}
\bibliography{references}

\end{document}
